\section{Prescribed edge and subgraph densities}\label{sec:fixed-densities}

In this section we prove Theorem~\ref{thm:fixed-densities} and
Corollary~\ref{cor:almost-every-tau}.  Fix a graph $H$ with $\Delta(H)\ge2$
and $\eps\in(0,1)$, and recall
$\tau_{\max}(\eps)=\max\{t(H,W):e(W)=\eps\}$.  By~\eqref{eq:entropy-identity},
the maximizers of~\eqref{eq:entropy-maximization} are exactly the minimizers of
\begin{equation}\label{eq:fixed-problem}
  \Psi_\eps(\tau):=\min\bigl\{I_\eps(W):e(W)=\eps,\ t(H,W)=\tau\bigr\},
  \qquad
  \eps^m<\tau<\tau_{\max}(\eps),
\end{equation}
and we work with this formulation.  The interval of values of $\tau$ is
nonempty: the graphon $\mathbf 1_{[0,\sqrt\eps]^2}$ has edge density $\eps$ and
$H$-density $\eps^{N/2}$, and $\eps^{N/2}>\eps^m$ because $2m\ge N+1$ when
$\Delta(H)\ge2$.  Every value in the interval is attained along the segment
$\eps+\theta(W_{\max}-\eps)$, $\theta\in[0,1]$, where $W_{\max}$ attains
$\tau_{\max}(\eps)$, and minimizers exist by compactness.

\subsection{The constraint qualification}

Inserting a block of vertices changes both the edge density and the
$H$-density, and both must be restored.  This requires two independent
directions of perturbation.

\begin{defn}[Constraint qualification]\label{def:constraint-qualification}
Let $W$ be a graphon with $e(W)=\eps$ and $t(H,W)=\tau$.  For $\eta>0$, put
$\Omega_\eta(W):=\{(x,y):\eta<W(x,y)<1-\eta\}$.  We say that $W$
\emph{satisfies the constraint qualification} if there are $\eta>0$ and bounded
symmetric functions $U_1,U_2$ supported in $\Omega_\eta(W)$ such that the matrix
\begin{equation}\label{eq:cq-matrix}
  \Gamma_W(U_1,U_2):=
  \begin{pmatrix}
    \int U_1 & \int U_2\\[1mm]
    \int\mathcal D_WU_1 & \int\mathcal D_WU_2
  \end{pmatrix}
\end{equation}
is invertible.  Otherwise we call $W$ \emph{degenerate}.
\end{defn}

The rows of $\Gamma_W(U_1,U_2)$ are the first variations of $e$ and
$t(H,\cdot)$ at $W$ in the directions $U_1$ and $U_2$.  Thus the constraint
qualification says that the derivative of $V\mapsto(e(V),t(H,V))$ at $W$ is
surjective onto $\RR^2$, as required in Lemma~\ref{lem:lagrange}, even when
restricted to bounded perturbations supported where $W$ is bounded away from
$0$ and $1$.  The support condition ensures that small corrections along
$U_1,U_2$ keep the graphon inside $[0,1]$.

\subsection{Monotonicity and positive multipliers}

\begin{lemma}[Monotonicity]\label{lem:fixed-monotone}
The function $\tau\mapsto\Psi_\eps(\tau)$ is strictly increasing on
$(\eps^m,\tau_{\max}(\eps))$.  Moreover, a minimizer of~\eqref{eq:fixed-problem}
at $\tau$ also minimizes $I_\eps$ over all graphons $U$ with $e(U)=\eps$ and
$t(H,U)\ge\tau$.
\end{lemma}

\begin{proof}
Let $\eps^m<\tau_1<\tau_2<\tau_{\max}(\eps)$, and let $W_2$ be a minimizer at
$\tau_2$.  The graphons $W_\theta:=\eps+\theta(W_2-\eps)$ have edge density
$\eps$, and $\theta\mapsto t(H,W_\theta)$ is continuous with values $\eps^m$ and
$\tau_2$ at $\theta=0,1$; hence $t(H,W_{\theta_0})=\tau_1$ for some
$\theta_0\in(0,1)$.  By strict convexity of $J_\eps$ and $J_\eps(\eps)=0$,
$\Psi_\eps(\tau_1)\le I_\eps(W_{\theta_0})<\theta_0I_\eps(W_2)\le\Psi_\eps(\tau_2)$.
For the second claim, if $e(U)=\eps$ and $t(H,U)>\tau$, the same interpolation
between $\eps$ and $U$ produces a graphon with $H$-density $\tau$ and cost
strictly less than $I_\eps(U)$; hence $I_\eps(U)>\Psi_\eps(\tau)$.
\end{proof}

\begin{lemma}[Interiority]\label{lem:fixed-interior}
Every minimizer $W$ of~\eqref{eq:fixed-problem} that satisfies the constraint
qualification satisfies $\gamma\le W\le1-\gamma$ for some $\gamma>0$.
\end{lemma}

\begin{proof}
Fix $\eta$, $U_1$, $U_2$ as in Definition~\ref{def:constraint-qualification}, and
write $\mathbf C(V):=(e(V),t(H,V))$.  Suppose that $W$ is not essentially
bounded away from $0$.  For large $\Lambda>0$ with $e^{-\Lambda}<\eta$, the symmetric set
$E:=\{W<e^{-\Lambda}\}$ has positive measure and is disjoint from the supports of
$U_1,U_2$.  Choose $(c_1,c_2)$ with
$\Gamma_W(U_1,U_2)(c_1,c_2)^\top=-D\mathbf C(W)[\mathbf 1_E]$.  Since
$0\le\mathcal D_W\le m$, we have $|c_1|+|c_2|\le C|E|$ with $C$ independent of
$\Lambda$, and $Z:=\mathbf 1_E+c_1U_1+c_2U_2$ satisfies $D\mathbf C(W)[Z]=0$.  The map
$(\theta,\xi_1,\xi_2)\mapsto\mathbf C(W+\theta Z+\xi_1U_1+\xi_2U_2)$ is
polynomial with invertible derivative in $(\xi_1,\xi_2)$ at the origin, so the
implicit function theorem gives $\xi_i(\theta)=O_\Lambda(\theta^2)$ with
$\mathbf C(V_\theta)=\mathbf C(W)$ for
$V_\theta:=W+\theta Z+\xi_1(\theta)U_1+\xi_2(\theta)U_2$.  For small $\theta>0$,
$V_\theta$ is a graphon.  If $0\le z<e^{-\Lambda}$ and $0<\theta<e^{-\Lambda}$, then
$J_\eps(z+\theta)-J_\eps(z)\le-\theta(\Lambda-C_\eps)$ for a constant $C_\eps$
independent of $\Lambda$, while $J_\eps'$ is bounded on $[\eta/2,1-\eta/2]$.
Therefore
\[
  I_\eps(V_\theta)-I_\eps(W)\le-\theta(\Lambda-C_\eps)|E|+C'\theta(|c_1|+|c_2|)+O_\Lambda(\theta^2)
  \le-\theta(\Lambda-C'')|E|+O_\Lambda(\theta^2),
\]
which is negative for $\Lambda$ large and then $\theta>0$ small, contradicting
minimality.  The case of values near $1$ is analogous, with $E:=\{W>1-e^{-\Lambda}\}$
and $-\mathbf 1_E$ in place of $\mathbf1_E$.
\end{proof}

\begin{lemma}[Positive multiplier]\label{lem:fixed-multiplier}
Let $W$ be a minimizer of~\eqref{eq:fixed-problem} that satisfies the
constraint qualification.  There are $\alpha\in\RR$ and $\mu>0$ such that
\begin{equation}\label{eq:fixed-el}
  J_\eps'\bigl(W(x,y)\bigr)=\alpha+\mu\,\mathcal D_W(x,y)
  \qquad\text{for almost every }(x,y).
\end{equation}
\end{lemma}

\begin{proof}
By Lemma~\ref{lem:fixed-interior}, $I_\eps$ is continuously differentiable on a
small $L^\infty$-ball around $W$ consisting of graphons, and the constraint
qualification says that $D\mathbf C(W)$ is surjective onto $\RR^2$.
Lemma~\ref{lem:lagrange} gives $\alpha,\mu\in\RR$ with
$DI_\eps(W)=\alpha\,De(W)+\mu\,Dt(H,W)$, and~\eqref{eq:fixed-el} follows as in
the proof of Lemma~\ref{lem:upper-tail-multiplier}.  Since the rows of
$\Gamma_W(U_1,U_2)$ are linearly independent, some linear combination $V$ of
$U_1,U_2$ satisfies $\int V=0$ and $\int\mathcal D_WV>0$.  Then $e(W+\theta V)=\eps$
and $t(H,W+\theta V)>\tau$ for small $\theta>0$, so
Lemma~\ref{lem:fixed-monotone} gives $DI_\eps(W)[V]\ge0$, that is,
$\mu\int\mathcal D_WV\ge0$ and $\mu\ge0$.  If $\mu=0$, then $J_\eps'(W)$ is
constant, so $W$ is constant; the edge constraint forces $W\equiv\eps$, which
contradicts $\tau>\eps^m$.
\end{proof}

\subsection{Inserting a block of vertices}

\begin{lemma}[Insertion inequality]\label{lem:fixed-insertion}
Let $W$ be a minimizer of~\eqref{eq:fixed-problem} that satisfies the
constraint qualification, and let $\alpha,\mu$ be as in
Lemma~\ref{lem:fixed-multiplier}.  Let $\mathcal F_W$ be the row
functional~\eqref{eq:row-functional} with $q=\eps$.  Then
\[
  \mathcal F_W(u)\ge c_W:=2I_\eps(W)-2\alpha\eps-\mu N\tau
  \qquad\text{for every measurable }u:[0,1]\to[0,1],
\]
and $\mathcal F_W(W_x)=c_W$ for almost every $x\in[0,1]$.
\end{lemma}

\begin{proof}
Fix $u$, and insert a block of measure $\theta$ with row $u$ as in the proof of
Lemma~\ref{lem:upper-tail-insertion}, with the value $\eps$ on the new block
and with the old values replaced by $W+\xi_1U_1+\xi_2U_2$, where $U_1,U_2$ are
as in Definition~\ref{def:constraint-qualification}.  Call the result
$W_{\theta,\xi}$.  Its edge density and $H$-density are polynomials in
$(\theta,\xi_1,\xi_2)$, with
\[
  \partial_\theta e(W_{\theta,0})\big|_{\theta=0}=2\Bigl(\int_0^1u-\eps\Bigr),
  \qquad
  \partial_\theta t(H,W_{\theta,0})\big|_{\theta=0}=\mathcal R_W(u)-N\tau,
\]
and their derivative in $(\xi_1,\xi_2)$ at the origin is
$\Gamma_W(U_1,U_2)$.  The implicit function theorem gives
$\xi(\theta)$ with $\xi(0)=0$ such that $W_\theta:=W_{\theta,\xi(\theta)}$
satisfies $e(W_\theta)=\eps$ and $t(H,W_\theta)=\tau$, and
\[
  De(W)[\xi_1'(0)U_1+\xi_2'(0)U_2]=-2\Bigl(\int_0^1u-\eps\Bigr),
  \qquad
  Dt(H,W)[\xi_1'(0)U_1+\xi_2'(0)U_2]=-\bigl(\mathcal R_W(u)-N\tau\bigr).
\]
For small $\theta>0$, $W_\theta$ is a graphon because the $U_i$ are supported
where $\eta<W<1-\eta$.  As before,
\[
  I_\eps(W_\theta)=(1-\theta)^2\int J_\eps\bigl(W+\xi_1(\theta)U_1+\xi_2(\theta)U_2\bigr)
  +2\theta(1-\theta)\int_0^1J_\eps(u)+\theta^2J_\eps(\eps),
\]
and using $DI_\eps(W)=\alpha\,De(W)+\mu\,Dt(H,W)$ we obtain
\[
  0\le\frac{\dd}{\dd\theta}\Big|_{\theta=0^+}I_\eps(W_\theta)
  =2\int_0^1J_\eps(u)-2I_\eps(W)-2\alpha\Bigl(\int_0^1u-\eps\Bigr)
  -\mu\bigl(\mathcal R_W(u)-N\tau\bigr)
  =\mathcal F_W(u)-c_W.
\]
Finally, $\int_0^1\int_0^1J_\eps(W_x(y))\dd y\dd x=I_\eps(W)$,
$\int_0^1\int_0^1W_x(y)\dd y\dd x=\eps$, and~\eqref{eq:rooted-average} give
$\int_0^1\mathcal F_W(W_x)\dd x=c_W$, so equality holds for almost every $x$.
\end{proof}

\begin{proof}[Proof of Theorem~\ref{thm:fixed-densities}]
Let $W$ be a maximizer of~\eqref{eq:entropy-maximization} that satisfies the
constraint qualification; equivalently, a minimizer of~\eqref{eq:fixed-problem}.
Apply Theorem~\ref{thm:row-criterion} with $q=\eps$, $\omega=\gamma$ from
Lemma~\ref{lem:fixed-interior}, and $\alpha,\mu$ from
Lemma~\ref{lem:fixed-multiplier}; its hypotheses are provided by
Lemma~\ref{lem:fixed-insertion}.
\end{proof}

\subsection{Degenerate maximizers occur only for a null set of values}

\begin{lemma}\label{lem:degenerate-forms}
Let $W$ be a degenerate minimizer of~\eqref{eq:fixed-problem}.  Then either
the set $\{W=0\}\cup\{W=1\}$ has positive measure, or $0<W<1$ almost everywhere
and $\mathcal D_W$ is almost everywhere equal to a constant.
\end{lemma}

\begin{proof}
Suppose that $0<W<1$ almost everywhere, and let $\eta>0$ be such that
$A:=\Omega_\eta(W)$ has positive measure.  With $\bar D:=|A|^{-1}\int_A\mathcal D_W$,
$U_1:=\mathbf 1_A$, and $U_2:=(\mathcal D_W-\bar D)\mathbf 1_A$, we have
$\det\Gamma_W(U_1,U_2)=|A|\int_A(\mathcal D_W-\bar D)^2$.  Degeneracy forces
this determinant to vanish, so $\mathcal D_W$ is constant on $\Omega_\eta(W)$.
Letting $\eta\downarrow0$, $\mathcal D_W$ is constant almost everywhere.
\end{proof}

\begin{lemma}\label{lem:degenerate-slope}
If~\eqref{eq:fixed-problem} has a degenerate minimizer at
$\tau\in(\eps^m,\tau_{\max}(\eps))$, then
\[
  \limsup_{\tau'\uparrow\tau}\frac{\Psi_\eps(\tau)-\Psi_\eps(\tau')}{\tau-\tau'}=\infty.
\]
\end{lemma}

\begin{proof}
Let $W$ be a degenerate minimizer at $\tau$ and put $W_\theta:=(1-\theta)W+\theta\eps$
for $\theta\in[0,1]$.  Then $e(W_\theta)=\eps$ and $I_\eps(W_\theta)<I_\eps(W)$
for $\theta\in(0,1)$, so the second claim of Lemma~\ref{lem:fixed-monotone} gives
$\tau_\theta:=t(H,W_\theta)<\tau$.  Since $\tau_\theta\to\tau$ as
$\theta\downarrow0$,
\[
  \frac{\Psi_\eps(\tau)-\Psi_\eps(\tau_\theta)}{\tau-\tau_\theta}
  \ge\frac{I_\eps(W)-I_\eps(W_\theta)}{\tau-\tau_\theta},
\]
and it suffices to show that the right-hand side tends to infinity along
$\theta\downarrow0$.  By convexity,
$J_\eps((1-\theta)z+\theta\eps)\le(1-\theta)J_\eps(z)$ for every $z$.

If $\{W=0\}\cup\{W=1\}$ has positive measure, then, since
$J_\eps(0)-J_\eps(\theta\eps)$ and $J_\eps(1)-J_\eps(1-\theta(1-\eps))$ are both at
least $c\,\theta\log(1/\theta)$ for small $\theta$ and some $c>0$, we get
$I_\eps(W)-I_\eps(W_\theta)\ge c'\theta\log(1/\theta)$, while
$\tau-\tau_\theta\le C\theta$ because $\theta\mapsto\tau_\theta$ is a polynomial.
Otherwise, by Lemma~\ref{lem:degenerate-forms}, $\mathcal D_W$ is a constant
$c_0$, so
$\frac{\dd}{\dd\theta}\tau_\theta\big|_{\theta=0}=\int\mathcal D_W(\eps-W)=c_0(\eps-e(W))=0$,
and hence $\tau-\tau_\theta\le C\theta^2$, while
$I_\eps(W)-I_\eps(W_\theta)\ge\theta I_\eps(W)$ with $I_\eps(W)>0$.  In both cases
the ratio tends to infinity.
\end{proof}

\begin{proof}[Proof of Corollary~\ref{cor:almost-every-tau}]
By Lemma~\ref{lem:fixed-monotone}, $\Psi_\eps$ is monotone on
$(\eps^m,\tau_{\max}(\eps))$, so by Lebesgue's theorem it is differentiable with a
finite derivative at almost every point.  At such a point, the left difference
quotients converge to the finite derivative, so by
Lemma~\ref{lem:degenerate-slope} no degenerate minimizer exists there.  Hence
the set of values $\tau$ admitting a degenerate maximizer has measure zero, and
for every other $\tau$ Theorem~\ref{thm:fixed-densities} applies to every
maximizer.
\end{proof}
